%% The following is a directive for TeXShop to indicate the main file
%%!TEX root = diss.tex

%================================
\chapter{Further Work}
\label{ch:Further}
%================================

%================================
\section{Coefficient Extraction Equations}
%================================

The work so far is strongly hinting towards a theory of ordinary and partial coefficient extraction equations - OCEs and PCEs. Here we will work out some of the theory. We start with a partial coefficient extraction operator of finite type
\begin{defn}$\mc L$ is a coefficient operator if
\begin{align}
	\mc L = \sum_{\substack{I\subseteq [d]\\ J \subseteq \ze^d}}a_{I,J}[x_I^J] = K + \mc E\qandq \mc E  = \sum_{\substack{I\subseteq [d],I\neq \emptyset\\ J \subseteq \ze^d}}a_{I,J}[x_I^J]\,,
\end{align}
where $K$ is the kernel of the operator and $\mc E$ is the pure extraction part. All the coefficients are Laurent polynomials in the $x$ variables and perhaps some other variables. The operator is ordinary if $d = 1$ and partial if $d \ge 2$. The operator is either of monomial type, where the coefficients are monomials, or of polynomial type, if one or more coefficients is a Laurent polynomial in $x_1,\ldots,x_d$.
\end{defn}
\begin{defn} An ordinary/partial coefficient equation (OCE/PCE) is one of the form $\mc L Q = q$ for some functions $Q$ and $q$, where $\mc L$ is an ordinary/partial coefficient operator. The equation is homogeneous if $q = 0$ and inhomogeneous if $q \neq 0$.
\end{defn}
The question is whether we can solve these kinds of equations in general, and what their solutions may look like. In the previous section we found that if $\mc L$ is finite then the function $Q$ is D-finite in the $x$ variables while it may not necessarily be D-finite in the $t$ variable or any other variable (parameter) that appears in the coefficients.

\begin{exa} The operator $\Gamma$ is a partial coefficient operator with monomial coefficients.
\end{exa}

\begin{exa} We can transform the functional equation for a finite random walk model into a PCE. To do so, we note that $K$ is the same in both forms. The part requiring some work is the one that has to do with the term $(1-\Gamma)SQ$ which is a pure extraction of $SQ$. We can expand $S$ into its monomials. Then $(1-\Gamma)S$ fully expands into monomials of the form
\begin{align}
	a_{I,J}\ext{x_I^J}x_A^B\,,
\end{align}
where we can pull terms inside $[x_I^J]$ from the right. In the end the $S$ is absorbed within the operator and istead of having monomial coefficients it now has polynomial coefficients. It may be useful to write an operator $\mc L$ in the form
\begin{align}
	\mc L = K + \Gamma'S
\end{align}
for some monomial operator $\Gamma'$.
\end{exa}

%================================
\section{Connection With Fredholm Theory}
%================================
Suppose that we are given the equation in the form
\begin{gather}
	\begin{align}
		KQ +t\Gamma'SQ = q\,,
	\end{align}
	\shortintertext{where}
	\begin{align}
		\Gamma' = \sum_{I,J}a_{I,J} x_I^J[x_I^J]\,,
	\end{align}
\end{gather}
where we only have finite-sized $J$'s and $\emptyset \neq I\subseteq [d]$. From complex analysis, we know that for a function $f(x)$ given by a power series and for $x$ within the region of convergence,
\begin{gather}
	\begin{align}
		[x^n]f(x) = f_n = \frac{1}{2\pi i}\int_\gamma \frac{f(z)}{z^{n+1}}dz
	\end{align}
	\shortintertext{and}
	\begin{align}
		f(x) = \frac{1}{2\pi i}\int_\gamma \frac{f(z)}{z-x}dz\,,
	\end{align}
\end{gather}
where in the first case, $\gamma$ winds CCW around the origin within the region of convergence, and in the second case, $\gamma$ winds around $x$ within the region of convergence. The idea is to combine these two expressions for many dimensions, at least formally, to represent the operator $\mc L$ in integral form. We have for $|I| = n$
\begin{align}
	[x_I^J]f(x_1,\ldots,x_d)= \frac{1}{(2\pi i)^{n}}\int\cdots\int_\gamma \frac{f}{z_I^{J+1}}dz_I\,,
\end{align}
where $f$ has $x_i = z_i$ for $i \in I$ and the rest are still $x$'s. If we use this representation we will have different numbers of iterations of the integral for different coefficient. We might as well go all-out and write
\begin{align}
	[x_I^J]f(x_1,\ldots,x_d)= \frac{1}{(2\pi i)^d}\int\cdots\int_\gamma \frac{f}{z_I^{J+1}(z-x)_{I^c}}dz_{[d]}\,,
\end{align}
where now $x_i = z_i$ in $f$ and 
\begin{align}
	(z-x)_P = \prod_{p\in P} (z_p-x_p)\,.
\end{align}
Therefore we may write the operator as
\begin{align}
	(K+t\Gamma'S)Q&= \frac{1}{(2\pi i)^d}\int\cdots\int_\gamma \left(\frac{K}{(z-x)_{[d]}} +t S\sum_{I,J}\frac{a_{I,J}}{z_I^{J+1}(z-x)_{I^c}}\right)Q dz_{[d]}\\
	& = \frac{1}{(2\pi i)^d}\int\cdots\int_\gamma \frac{1}{(z-x)_{[d]}}\left(K +t S\sum_{I,J}\frac{a_{I,J}(z-x)_I}{z_I^{J+1}}\right)Q dz_{[d]}\\
	&= \frac{1}{(2\pi i)^d}\int\cdots\int_\gamma R_0Q dz_{[d]}\,.
\end{align}
Now to turn this into a $d$-dimensional Fredholm problem we will need to parametrize the integral. One way that might work is to set $z_k = r_k \exp(i\theta_k)$ and $x_k = s_k e^{i\phi_k}$. Since $dz_k = r_k i e^{i\theta_k}$, we have
\begin{gather}
	\begin{align}
		(K+t\Gamma'S)Q = \frac{r_{[d]}}{(2\pi)^d}\int_0^{2\pi}\cdots\int_0^{2\pi} e^{i\theta_1}\ldots e^{i\theta_d}R_0Q d\theta_{[d]} = q(\phi_1,\ldots,\phi_d)\,,
	\end{align}
	\shortintertext{so if we write}
	\begin{align}
		R = \exp(i(\theta_1+\ldots + \theta_d))R_0(\phi_1,\ldots,\phi_d,\theta_1,\ldots,\theta_d)\,,
	\end{align}
\end{gather}
we have the inhomogeneous Fredholm equation of the first kind in $d$ dimensions
\begin{align}
	\int_0^{2\pi}\cdots\int_0^{2\pi} RQ d\theta_1\ldots d\theta_d= g\,,
\end{align}
where $R$ is the kernel and $g = (2\pi)^dq/r_{[d]}$ and $Q$ is in terms of $\theta_k$'s. Perhaps Fredholm theory can be applied in solving the random walk models in man dimensions. Perhaps this will also lead to a good inner product on the space where $Q$ lives.


\endinput
